\section{Sparse pinning and descent}
\label{sec:pinning}

\subsection{Sparse pinning and removal of unassigned vertices}

We retain the notation and hypotheses of Theorem~\ref{thm:measured-paving}.
In particular, $b>0$ and $r$ satisfy its strict largest-root inequality,
and $K,C$ are the universal constants from
Lemma~\ref{lem:buffered-pin}.  The lifted graph on a slab is always the
restriction of the original enclosing graph, including edges at which a
particular matrix entry happens to vanish.  All distances below refer to
this enclosing graph.

\begin{proposition}\label{prop:sparse-pinning}
For every $\zeta,\eta>0$ and $L_0\geq1$, there are an integer
$L\geq L_0$ and a measurable partition
\[
 V=\mathcal U\sqcup V_1\sqcup\cdots\sqcup V_r
\]
such that $\mu(\mathcal U)<\eta$ and the following holds.  On the
surviving lifted vertices $(V\setminus\mathcal U)\times[0,L]$, let
$T_{\mathrm{col}}^L$ be the restriction of $T$, keeping only entries
between vertices whose basepoints belong to the same $V_i$, and put
\begin{equation}\label{pin:surviving-operator}
 A_L=d^{1/2}(T_{\mathrm{col}}^L-Kb\one)d^{1/2},
 \qquad d_{(v,s)}=e^s.
\end{equation}
Extend $m_u^{\mathrm{surv}}=((A_L)_+)_{uu}$ by zero on the deleted
vertices.  Then
\begin{equation}\label{pin:surviving-cost}
 \frac1L\int_{V\times[0,L]}m_u^{\mathrm{surv}}\,
 d\widetilde\mu(u)<\zeta.
\end{equation}
\end{proposition}

\begin{proof}
We first construct a partial assignment with small mixed-polynomial
cost. Polynomial moments depend on finite graph balls, so sufficiently
sparse simultaneous pins have an additive first-order response except
on an event of quadratic order in the update probability. Approximating
the derivative of the positive-part function controls the error in
that response independently of the number of roots.
Parameters used to select the slab will be specified below.
On a fixed slab $[0,L]$, write
\[
 \Phi=\frac1L\int_{V\times[0,L]}m_u\,d\widetilde\mu(u).
\]
Choose a common spectral support bound $[-B,B]$ for all the finite
principal systems, pin alternatives, and deletion configurations on this
slab, as in Lemma~\ref{lem:slab-limits}.  For $\tau>0$, choose a real
polynomial $P$ such that
\begin{equation}\label{pin:polynomial-approximation}
 P(-B)=0,\qquad
 \int_{-B}^B\abs{P'(x)-\mathbf1_{(0,\infty)}(x)}\,dx\leq\tau.
\end{equation}
For example, approximate the indicated step function first by a
continuous function in $L^1([-B,B])$, then uniformly by a polynomial,
and integrate that polynomial from $-B$.  It follows that
\[
 \sup_{[-B,B]}\abs{P(x)-x_+}\leq\tau.
\]
Define
\[
 m_{P,u}=\int P(x)\,d\gamma_u(x),\qquad
 \Phi_P=\frac1L\int_{V\times[0,L]}m_{P,u}\,d\widetilde\mu(u).
\]
Since $\gamma_u$ is a probability measure and
$\widetilde\mu(V\times[0,L])=1-e^{-L}\leq1$, we have
\begin{equation}\label{pin:cost-approximation}
 \abs{\Phi-\Phi_P}\leq\tau \qquad(L\geq1).
\end{equation}

\paragraph{The response to one buffered pin.}
At an unassigned lifted vertex $u$, consider pinning to a specified color
$i$, increasing its penalty from $bd_u$ to $Kbd_u$ at the same time.
The total polynomial response on its component is
\begin{equation}\label{pin:single-response}
 \begin{split}
 D_P(u)&=\sum_y\bigl(m_{P,y}^{\mathrm{pin}\ u\mathrm{\ to\ }i}
                         -m_{P,y}^{\mathrm{pre}}\bigr)\\
 &=\int_{-B}^B P'(x)
       \bigl(\xi_{h_{u,i}'}(x)-\xi_{h_u}(x)\bigr)\,dx.
 \end{split}
\end{equation}
Here $h_{u,i}'$ denotes the buffered alternative; the prime does not
denote differentiation.  To prove the identity on a finite principal
set, use the fact that the sum of its local spectral probability
measures is the counting measure of the roots.  The pre-pin and
post-pin polynomials have equal degree and the same principal
polynomial with $u$ removed.  Integrating the difference of their root
counts above $x$ against $P'(x)$ gives exactly
\eqref{pin:single-response}.

Put $k=\max\{1,\deg P\}$.  By the locality in
Lemma~\ref{lem:slab-limits}, only vertices at distance at most $k$ from
$u$ contribute to the sum on the first line of
\eqref{pin:single-response}.  Compute in finite balls large enough to
contain these vertices and their relevant neighborhoods.  The local
polynomial moments stabilize, while the functions $\xi$ converge
weak-* on $[-B,B]$.  This proves the identity on infinite components
as well.  Since both $\xi$-functions take values in $[0,1]$ almost
everywhere, \eqref{pin:polynomial-approximation} implies
\begin{equation}\label{pin:single-error}
 \abs{D_P(u)-(\ell_{u,i}'-\ell_u)}\leq\tau.
\end{equation}

\paragraph{One sparse update.}
For a current partial base assignment, choose the prospective colors
$i(v)$ supplied by Lemma~\ref{lem:coherent-choice}, with tolerance
$\tau_0>0$.  The color $i(v)$ is the same at all heights over $v$.
Properly color the distance-$2k$ power of the base graph with finitely
many auxiliary Borel color classes.  The finite Borel coloring fact
used in the proof of Lemma~\ref{lem:far-deletion} applies because this
power still has bounded degree.  Independently for these auxiliary
classes, select a Bernoulli-$\alpha$ indicator, where $0<\alpha<1$.
Every currently unassigned base vertex in a selected class is pinned to
its prospective color, with its penalty buffered, at all heights
simultaneously.

Within a radius-$k$ lifted ball, different vertices project to different
base vertices, and their pairwise base distance is at most $2k$.
They consequently have different auxiliary classes.  The indicators of
all candidate pins affecting the local moment $m_{P,y}$ are therefore
independent.  This assertion uses injectivity of projection within
each lifted orbit; it does not require independent indicators at all
vertices of the base.

If at most one candidate in this local ball is selected, the change in
$m_{P,y}$ is exactly the sum of its selected single-pin responses.  A
ball has at most a fixed finite number $N$ of vertices, so the
probability of two or more selections is at most
$\binom{N}{2}\alpha^2$.  On this event the discrepancy between the
actual change and the sum of single-pin responses is bounded in
absolute value by a constant depending only on $N$ and
$\sup_{[-B,B]}\abs{P}$.  Integrate this estimate and exchange roots
and pin locations using \eqref{slab:mass-transport}.  The result is
\begin{equation}\label{pin:sparse-update}
 \E[\Phi_P^{\mathrm{new}}\mid\mathrm{pre}]
 \leq\Phi_P+
   \frac\alpha L\int_{\mathrm{unassigned}}D_P(u)\,
                  d\widetilde\mu(u)+C_P\alpha^2.
\end{equation}
All the transports here are measurable, uniformly bounded on the
fixed slab, and supported on bounded-radius pairs.  Thus they are
integrable and the signed version of mass transport is applicable.
The constant $C_P$ may depend on the slab, graph and polynomial, but
not on the current assignment, the prospective choices, or $\alpha$.

By Lemma~\ref{lem:coherent-choice},
\eqref{pin:single-error}, and \eqref{pin:cost-approximation},
\eqref{pin:sparse-update} gives
\begin{equation}\label{pin:recurrence}
 \E[\Phi_P^{\mathrm{new}}\mid\mathrm{pre}]
 \leq(1+C\alpha)\Phi_P+
       \alpha\bigl(\tau_0+(1+C)\tau+C_P\alpha\bigr).
\end{equation}
The same $L$ from Lemma~\ref{lem:coherent-choice} works for all partial
base assignments, so this step can be repeated.

\paragraph{Iteration and choice of parameters.}
Start with every vertex unassigned.  The initial cost is $\Phi=0$ by
Lemma~\ref{lem:slab-limits}.  Use fresh independent class indicators at
each step and recalculate the prospective colors from the current
assignment.  Given $S>0$, perform $n=\lceil S/\alpha\rceil$ steps.
For these fixed parameters there are only finitely many possible
histories.  At every node one may choose separately the measurable
prospective-color functions furnished by
Lemma~\ref{lem:coherent-choice}; no selection measurable in an
infinite-dimensional space of assignments is needed.

We have $n\alpha\leq S+1$ and
$(1+C\alpha)^n\leq e^{C(S+1)}$.  Iterating
\eqref{pin:recurrence}, using the initial bound $\Phi_P\leq\tau$,
and then using \eqref{pin:cost-approximation} again yields
\begin{equation}\label{pin:iteration}
 \E\Phi_{\mathrm{final}}
 \leq\tau+e^{C(S+1)}
   \bigl[\tau+(S+1)(\tau_0+(1+C)\tau+C_P\alpha)\bigr].
\end{equation}
If $\mathcal U$ is the final unassigned base set, each still-unassigned
vertex has conditional probability $\alpha$ of being selected at the
next step, regardless of its prospective color.  Therefore
\begin{equation}\label{pin:unassigned-measure}
 \E\mu(\mathcal U)=(1-\alpha)^n\leq e^{-S}.
\end{equation}

Here is an order of choices which makes both expectations small.
First take $S$ so large that $e^{-S}<\eta/4$.  Next take
$\tau_0,\tau>0$ so small that the terms in
\eqref{pin:iteration} not involving $C_P\alpha$ have sum less than
$\zeta/8$.  Choose an integer $L\geq L_0$ using
Lemma~\ref{lem:coherent-choice} for this $\tau_0$.  Now choose $B$ and
$P$ as above, thereby fixing $C_P$.  Finally take $\alpha$ small
enough that the remaining term in \eqref{pin:iteration} is less than
$\zeta/8$.  Thus
\[
 \E\Phi_{\mathrm{final}}<\zeta/4,
 \qquad \E\mu(\mathcal U)<\eta/4.
\]
Both random variables are nonnegative.  By Markov's inequality and a
union bound, some history satisfies simultaneously
\begin{equation}\label{pin:chosen-history}
 \Phi_{\mathrm{final}}<\zeta,
 \qquad\mu(\mathcal U)<\eta.
\end{equation}
Fix such a history.

\paragraph{Removing the unassigned vertices.}
We show that deleting $\mathcal U\times[0,L]$ cannot increase the
integrated positive cost.  Keep the initial and final deletion states
fixed.  For an arbitrary $\vartheta>0$, choose another polynomial $Q$
with
\[
 Q(-B)=0,\qquad
 \int_{-B}^B\abs{Q'(x)-\mathbf1_{(0,\infty)}(x)}\,dx
 \leq\vartheta.
\]
At each intermediate deletion state, extend $m_{Q,y}$ by zero at
vertices that have already been deleted.

For one deletion at $u$, the total pre-minus-post polynomial cost is
\begin{equation}\label{pin:single-deletion}
 \sum_y\bigl(m_{Q,y}^{\mathrm{pre}}-m_{Q,y}^{\mathrm{post}}\bigr)
     =\int_{-B}^B Q'(x)\xi_{h_u}(x)\,dx\geq-\vartheta.
\end{equation}
For finite sets the identity again follows by integrating the
difference of root counts; the boundary term caused by the loss of one
root is $Q(-B)=0$.  The inequality uses $0\leq\xi_{h_u}\leq1$.
The finite-range locality argument used for
\eqref{pin:single-response} proves the identity in an infinite
component.

Partition the set to be deleted into finitely many Borel batches whose
distinct vertices have distance greater than
$2\max\{1,\deg Q\}$ within each batch in the enclosing lifted
graph.  For one batch, at most one deletion affects any local
polynomial moment, so the batch change is exactly the sum of
single-deletion responses in its pre-state.  Mass transport,
\eqref{pin:single-deletion}, and telescoping over the batches show
that the integrated pre-minus-post polynomial cost, divided by $L$,
is at least
\[
 -\vartheta\,
   \frac{\widetilde\mu(\mathcal U\times[0,L])}{L}.
\]
The error is charged once to each deleted vertex; it is not multiplied
by the number of batches.  Intermediate deleted sets need not be
constant in the height variable, since Lemma~\ref{lem:slab-limits}
allows arbitrary deletion configurations in the enclosing slab.

At each endpoint the integrated error between the $Q$-cost and the
positive cost is at most $\vartheta$.  Thus the pre-minus-post
positive cost is at least $-3\vartheta$.  The endpoints are independent
of $Q$ and of the batch decomposition, so letting $\vartheta\downarrow0$
proves the desired monotonicity.

All surviving vertices are now deterministically colored, with penalty
$Kbd_u$.  Their finite mixed polynomial is the characteristic
polynomial of the corresponding restriction of
\eqref{pin:surviving-operator}, so its local measure is the diagonal
spectral probability measure of that operator.  On an infinite
component the same assertion follows from stabilization of polynomial
moments and the common compact spectral support.  Hence the surviving
positive cost is exactly $m_u^{\mathrm{surv}}$.  Monotonicity and
\eqref{pin:chosen-history} give \eqref{pin:surviving-cost}.
\end{proof}

\subsection{Descent from a slab to the base relation}

The next estimate converts the integrated weighted cost into the desired
spectral error on the base. A finite ball with a large eigenvalue forces
a definite positive cost in its lift; mass transport then bounds the
measure of such bad balls. For an integer $k\geq1$, let $N_k$ be a
uniform upper bound for the number of vertices in a radius-$k$ ball of
the base graph.

\begin{lemma}\label{lem:descent}
Fix $\sigma,\beta>0$, and choose a real polynomial $p$ of degree at most
$k\geq1$ such that
\[
 \sup_{x\in[-1,1]}
 \abs{p(x)-(x-(Kb+\sigma))_+^2}\leq\beta.
\]
Suppose a base partial assignment, with unassigned set $\mathcal U$,
and a slab length $L\geq\max\{1,4k\log J\}$ have surviving cost
\[
 \Psi=\frac1L\int_{V\times[0,L]}m_u^{\mathrm{surv}}\,
                         d\widetilde\mu(u),
\]
where $m^{\mathrm{surv}}$ is defined from
\eqref{pin:surviving-operator}.  Complete the base coloring by giving
$\mathcal U$ one additional color, and denote its paved matrix system
by $T_{\mathrm{pav}}$.  Then
\begin{equation}\label{pin:descent-bound}
 \int_V\bigl((T_{\mathrm{pav}}-(Kb+\sigma)\one)_+^2\bigr)_{vv}
                   \,d\mu(v)
 \leq2\beta+\mu(\mathcal U)+\frac{2J^kN_k}{\sigma}\Psi.
\end{equation}
\end{lemma}

\begin{proof}
Put $g(x)=(x-(Kb+\sigma))_+^2$.  Since the paved operator is a
compression by mutually orthogonal color projections, it remains a
self-adjoint contraction; in particular $0\leq g(T_{\mathrm{pav}})
\leq\one$.

For an assigned base vertex $v$, let $B_k(v)$ be its radius-$k$ ball
in the base graph, and let $T_v$ be the finite compression of the paved
matrix to the assigned vertices of this ball.  Polynomial diagonal
moments through degree $k$ at $v$ are unchanged by this compression.
Indeed every walk contributing to such a moment stays within $B_k(v)$,
and edges from assigned to unassigned vertices are zero after paving.
Consequently
\[
 (p(T_{\mathrm{pav}}))_{vv}=(p(T_v))_{vv}.
\]
Let $\mathcal F$ be the measurable set of assigned vertices for which
$\lambda_{\max}(T_v)>Kb+\sigma$.  Measurability follows by finite
matrix evaluation on the Borel balls.  If $v\notin\mathcal F$ is
assigned, then $g(T_v)=0$ and the two polynomial approximation errors
give
\[
 (g(T_{\mathrm{pav}}))_{vv}\leq2\beta.
\]
On $\mathcal F\cup\mathcal U$ its diagonal is at most $1$.  Therefore
\begin{equation}\label{pin:failure-reduction}
 \int_V(g(T_{\mathrm{pav}}))_{vv}\,d\mu(v)
 \leq2\beta+\mu(\mathcal U)+\mu(\mathcal F).
\end{equation}

It remains to bound $\mu(\mathcal F)$.  Set
$I_{\mathrm{int}}=[k\log J,L-k\log J]$.
For $u=(v,s)$ with $s\in I_{\mathrm{int}}$, the entire ball $B_k(v)$
lifts into the slab: the height change along each graph edge has
absolute value at most $\log J$.  Let $Q_u$ project onto the assigned
vertices of this lifted ball.  With $A_L$ as in
\eqref{pin:surviving-operator}, the order inequality $A_L\leq(A_L)_+$
gives, after compression and diagonal congruence,
\begin{equation}\label{pin:positive-upper-bound}
 T_v-Kb\one
 \leq d_{\mathrm{ball}}^{-1/2}Q_u(A_L)_+Q_u
                                  d_{\mathrm{ball}}^{-1/2}.
\end{equation}
The base and lifted finite matrices here are identified up to the
allowed diagonal unitary gauge, which does not affect the inequality
or any diagonal spectral quantities.  The right side is positive.
If $v\in\mathcal F$, its norm is greater than $\sigma$ by
\eqref{pin:positive-upper-bound}, and hence its trace is greater than
$\sigma$.  Equivalently,
\begin{equation}\label{pin:failure-trace}
 \sum_{y\in B_k(u)\ \mathrm{assigned}}
                \frac{m_y^{\mathrm{surv}}}{d_y}>\sigma.
\end{equation}
Here $B_k(u)$ denotes the lift of the base ball, which agrees with the
radius-$k$ enclosing lifted ball at these interior heights.  The
comparison uses only $A_L\leq(A_L)_+$; it does not require positive
part to commute with compression or diagonal congruence.

Integrate \eqref{pin:failure-trace} against $d\mu(v)\,ds/L$ over
$\mathcal F\times I_{\mathrm{int}}$, and enlarge the resulting
nonnegative sum to all interior roots.  Since
$d\mu(v)\,ds=d_u\,d\widetilde\mu(u)$, the resulting integrand is a
sum of $(d_u/d_y)m_y^{\mathrm{surv}}$.  On a radius-$k$ ball,
$d_u/d_y\leq J^k$.  Exchanging roots and targets by
\eqref{slab:mass-transport}, each target $y$ has at most $N_k$
possible roots.  Deleted vertices may be retained as locations for
the enclosing-ball relation, with $m_y^{\mathrm{surv}}=0$ there.
Thus
\[
 \sigma\frac{\abs{I_{\mathrm{int}}}}L\mu(\mathcal F)
 \leq J^kN_k\Psi.
\]
Our lower bound on $L$ ensures
$\abs{I_{\mathrm{int}}}/L\geq1/2$, including the case $J=1$.
We conclude that
\[
 \mu(\mathcal F)\leq\frac{2J^kN_k}{\sigma}\Psi.
\]
Combining this with \eqref{pin:failure-reduction} proves
\eqref{pin:descent-bound}.
\end{proof}

\begin{proof}[Completion of the proof of Theorem~\ref{thm:measured-paving}]
Fix the prescribed $\sigma,\rho>0$.  Choose $\beta>0$ with
$2\beta<\rho/4$, and then choose the polynomial and its degree
$k\geq1$ in Lemma~\ref{lem:descent}.  Thus $N_k$ is fixed before
selecting the slab.  Take
\[
 \eta=\rho/4,\qquad
 \zeta=\frac{\rho\sigma}{8J^kN_k},\qquad
 L_0=\max\{1,4k\log J\}.
\]
Proposition~\ref{prop:sparse-pinning} supplies a partial assignment and
a slab of length at least $L_0$ with
$\mu(\mathcal U)<\eta$ and $\Psi<\zeta$.  Give $\mathcal U$ the
additional color and apply Lemma~\ref{lem:descent}.  Its right side is
strictly less than $\rho$, as required.

Only the $r$ assigned colors and one extra color occur in the final
partition.  The auxiliary graph colors, the slab length, the
approximation degrees and the number of updates affect its
construction but not this number of final colors.  In particular the
graph and cocycle bounds do not enter the condition on $r$ in the
theorem.
\end{proof}
